We present a sparse-view pipeline for constructing language-grounded 3D scene graphs
from casually captured RGB images. Feed-forward geometry models can now predict
cameras, depth, and world-point maps without per-scene optimization, while 2D
foundation models provide class-agnostic masks and open-vocabulary descriptors.
However, these outputs remain view-local: they do not directly identify persistent
objects, expose uncertain geometric evidence, or produce relation-bearing scene
structure. Our method bridges this gap by lifting SAM proposals into VGGT world-point
maps, attaching CLIP and DINOv2 descriptors, scoring lifted candidates with confidence
and compactness uncertainty, fusing cross-view candidates into object nodes, and
exporting open-vocabulary labels and spatial relations. The graph representation keeps
symbolic predictions together with source views, proposal support, 3D point samples,
and uncertainty values, making the output inspectable and reviewable. We evaluate the
pipeline on a five-sequence TUM RGB-D sparse-view benchmark with 20 runs across 3, 5,
8, and 10 input views. Averaged over scenes, moving from 3 to 10 views increases lifted
proposals from 54.0 to 173.6, fused objects from 31.8 to 72.2, multi-view objects from
11.8 to 28.2, and inferred relations from 585.2 to 2104.0. Against checked references
initialized from the 10-view graphs, non-reference object-label F1 improves from 0.622
to 0.914 and relation-triplet F1 from 0.434 to 0.799 between 3 and 8 views. These
reference-consistency scores measure sparse-view recovery rather than dense-scene
recognition accuracy, establishing a reproducible graph construction baseline and a
concrete path from predicted graphs to object-relation review.
